\section{Materiais e Métodos}
\label{sc:materiais_e_metodos}

\subsection{Padrão Arquitetural e Stack Tecnológico}
O sistema adota o padrão \textbf{FCIS (\textit{Functional Core, Imperative/OO Shell})}. As transformações matemáticas, sanitizações e cálculos de elipse são modelados como funções puras, sem estado e vetorizadas (\textit{DataFrame $\rightarrow$ DataFrame}). A casca orientada a objetos é restrita ao gerenciamento determinístico de conexões e soquetes (Kafka, Druid, FastMCP) via \textit{Context Managers} com injeção de dependência. O ecossistema é operado em Python 3.12 sob governança do gerenciador Astral \texttt{uv}.

A infraestrutura tecnológica integra: (i) Dagster para orquestração; (ii) Apache Kafka para mensageria; (iii) Apache Druid (OLAP colunar particionado por tempo) e DuckDB (\textit{spatial} e \textit{h3}) para persistência em GeoParquet; (iv) Apache Superset com camadas H3 (\textit{deck.gl}) e biblioteca \texttt{h3-py} para indexação em \texttt{uint64}; (v) Servidor FastMCP para interface LLM determinística, templates WeasyPrint/Markdown e despachante assíncrono de Webhooks; e (vi) Bases de dados: INPE VIIRS (375 m) e GOES-19 (10 min), reanálise horária ERA5, dados abertos ONS e MapBiomas Fogo.

\subsection{Procedimentos Metodológicos}
A metodologia estrutura-se em cinco etapas sequenciais:
\begin{enumerate}
    \setlength\itemsep{0em}
    \item \textbf{Ingestão Contratual (Bronze):} Coleta de anomalias térmicas e séries climáticas com validação via Pydantic;
    \item \textbf{Indexação H3 e Modelagem Anisotrópica (Silver):} Discretização das feições de interesse nas resoluções H3 7 ($\sim 5\text{ km}^2$) e 9 ($\sim 0,1\text{ km}^2$). Cálculo dos semieixos ($a, b$) e azimute da elipse acoplados ao vetor de vento e à carga combustível local;
    \item \textbf{Agregação e Carga Analítica (Gold):} Cruzamento das células H3 ativas contra a infraestrutura via \textit{hash joins} $O(1)$. Carga da tabela Fato no Druid no modelo analítico Snowflake (estritamente FKs numéricas, métricas e timestamps);
    \item \textbf{Inteligência Ativa (MCP e Webhooks):} Avaliação periódica de limiares críticos com máquina de \textit{cooldown} (janela de 6h). Ativação do FastMCP para executar \textit{golden queries} no Druid, disparando o boletim SitRep (ICS-209);
    \item \textbf{Validação Experimental:} Comparação de falsos positivos entre buffers circulares e elípticos, auditoria de latência de ponta a ponta ($< 30\text{ s}$) e teste de consistência factual das respostas do agente de IA.
\end{enumerate}
